\documentclass{beamer}

\usepackage[utf8]{inputenc}

\input{../helper}

\title{Introdução à Computação \\ selection sort}
\author{Alexandre Rademaker}
\date{}

\begin{document}

\begin{frame}
 \maketitle
\end{frame}


\begin{frame}[fragile,allowframebreaks]{Ordenação (sort)}

  Dado uma lista de valores, possívelmente desordenada, queremos
  escrever um algorítmo que ordena esta lista, segundo algum critério.

  Por exemplo, para a lista
  \[
    [ 6 \; 3 \; 8 \; 5 \; 2 \; 7 \; 4 \; 1 ]
  \]
  queremos produzir
  \[
    [ 1 \; 2 \; 3 \; 4 \; 5 \; 6 \; 7 \; 8 ]
  \]

  \framebreak

  Vamos considerar dois algorítmos para esta tarefa:

  \begin{itemize}
  \item \emph{selection sort}
  \item bubble sort
  \item merge sort
  \end{itemize}

\end{frame}


\begin{frame}[fragile,allowframebreaks]{Selection Sort}

  Dada uma lista de números

  \begin{lstlisting}[style=code,basicstyle=\ttfamily\scriptsize]
    For i from 0 to n - 1
     Find smallest between nums[i] and nums[n - 1]
     Swap smallest number with nums[i]
  \end{lstlisting}

  \framebreak

  \begin{lstlisting}[style=code]
    5 2 7 4 1 6 3 0

    5 2 7 4 1 6 3 0
                  ^

    0 | 2 7 4 1 6 3 5
    
    0 | 2 7 4 1 6 3 5
              ^    

    0 1 | 7 4 2 6 3 5
              ^
	  
    0 1 2 | 4 7 6 3 5
                  ^
  \end{lstlisting}

  \framebreak

  Para este algorítmo, considerando uma lista de $n$
  elementos. Começamos olhando para $n$ elementos, depois para $n - 1$
  etc.

  \begin{align*}
    n + (n - 1) + (n - 2) + \ldots + 1 & = \\
    \frac{(n^2 + n)}{2} = \frac{n^2}{2} + \frac{n}{2} & = O(n^2) 
  \end{align*}

  No melhor caso, quando a lista já estiver ordenada, ainda assim
  temos que olhar para todos os números, logo também $\Omega(n^2)$.
  
  \framebreak

  \begin{lstlisting}[style=normalc,caption=src/selection.c]
    void selection_sort (int a[], int n) {
      int m, t;
      for (int i = 0; i < n; i++) {
        for (int j = i, m = i; j < n; j++) {
          if (a[j] < a[m])
          m = j;
        }
        t = a[i];
        a[i] = a[m];
        a[m] = t;
      }
    }
  \end{lstlisting}

\end{frame}


\end{document}

%%% Local Variables:
%%% mode: latex
%%% TeX-master: t
%%% End:
